Table \ref{tab.readingList} collects the reading of the course.
All reading is optional and meant for the student who wants to go further.

\begin{table}[htbp]
 \centering
 \begin{tabular}{p{0.26\textwidth} p{0.63\textwidth}}
  \toprule
  Reference & What it is, and why it is on the list \\
  \midrule
  \multicolumn{2}{l}{\emph{Mathematical background}} \\
  \cite{Wil91pro} &
   Measure-theoretic probability, conditional expectation and discrete-time martingales.
   The background Section \ref{sec.prerequisites} assumes. \\
  \addlinespace
  \multicolumn{2}{l}{\emph{Python}} \\
  \cite{Ram22flu} &
   The data model, the protocols, the iterators and the type hints,
   in the order in which the language is built.
   What idiomatic Python is, and why. \\
  \cite{Sla24eff} &
   A list of choices, each with the reason for it and the cost of the alternative.
   The judgment of Section \ref{sec.theCourse}, exercised item by item. \\
  \addlinespace
  \multicolumn{2}{l}{\emph{Market microstructure}} \\
  \cite{CKS14pri} &
   Introduces the order flow imbalance and regresses price changes on it.
   The reference of Section \ref{sec.priceFormation}. \\
  \cite{FPR13mar} &
   The institutions of an order-driven market and the economics of liquidity.
   Chapters 1--3 cover the ground of Section \ref{sec.orderDrivenMarkets}. \\
  \cite{Bou18tra} &
   The empirical statistics of the limit order book, and the models fitted to them. \\
  \addlinespace
  \multicolumn{2}{l}{\emph{Option pricing and volatility trading}} \\
  \multicolumn{2}{l}{Published with Chapter \ref{chap.options}.} \\
  \bottomrule
 \end{tabular}
 \caption{The reading of the course.}
 \label{tab.readingList}
\end{table}
